\documentclass[11pt]{article}

% Packages
\usepackage{amsmath}
\usepackage{amssymb}
\usepackage{amsthm}
\usepackage{geometry}
\usepackage[colorlinks=true, linkcolor=blue, toccolor=blue]{hyperref}
\usepackage{mathrsfs}
\usepackage{mathtools}
\usepackage{tikz}

% Page Geometry
\geometry{letterpaper, margin=1in}

% Theorem Environments
\newtheorem{problem}{Problem}
\newtheorem{theorem}{Theorem}
\newtheorem{lemma}{Lemma}
\newtheorem{definition}{Definition}

% QED Symbol configuration
\renewcommand{\qedsymbol}{$\blacksquare$}

% Title Information
\title{\textbf{Advanced Practice Problem Set: \\ Rigorous Proofs Involving Cartesian Products of Sets}}
\author{Elite Mathematics Faculty}
\date{\today}

\begin{document}

\maketitle

\tableofcontents
\newpage

% ==========================================
% SECTION 1: CONCEPT CHECKLIST & MAPPING
% ==========================================
\section{Concept Checklist \& Question Mapping}
\label{sec:checklist_start}

The following theoretical concepts, logical structures, and proof techniques have been extracted from the source material and rigorously tested in this problem set.

\begin{itemize}
    \item \textbf{C1: Kuratowski Definition of Ordered Pairs} -- Formal set-theoretic definition of $(a,b)$. (Tested in P1)
    \item \textbf{C2: Power Sets \& Combinatorial Enumeration} -- Cardinality and power sets of Cartesian products. (Tested in P2, P3, P4, P5)
    \item \textbf{C3: Empty Sets \& Vacuous Truths} -- Handling $\emptyset$ as the "zero" of set theory; existential instantiation. (Tested in P6, P7, P8, P9, P10)
    \item \textbf{C4: Distributivity of Cartesian Products} -- Distributing over $\cup, \cap, \setminus$. (Tested in P11, P12, P13, P14, P15)
    \item \textbf{C5: Subset Inclusion \& Proof by Cases} -- Handling unions and logical "OR" statements. (Tested in P16, P17, P20)
    \item \textbf{C6: Set Equality via Logical Equivalence} -- Translating set operations to propositional logic ($\land, \lor$). (Tested in P14, P15, P18, P19)
    \item \textbf{C7: Complements of Cartesian Products} -- De Morgan's Laws applied to 2D spaces. (Tested in P18)
    \item \textbf{C8: Infinite Cartesian Products \& Axiom of Choice} -- Arbitrary intersections and non-empty families of sets. (Tested in P21, P22, P23)
    \item \textbf{C9: Abstract Real-World Modeling} -- Markov Decision Processes (MDPs), State-Action spaces, and Joint Probability. (Tested in P24, P25, P26, P27)
    \item \textbf{C10: Diagnostic Testing (Find the Flaw)} -- Identifying logical fallacies, type mismatches, and false assumptions. (Tested in P28, P29, P30)
\end{itemize}

\newpage

% ==========================================
% SECTION 2: PROBLEM SET
% ==========================================
\section{Problem Set}
\label{sec:problems}

\subsection{Fundamentals, Power Sets, and Combinatorics}
\begin{problem}
Using the Kuratowski definition of an ordered pair, $(a, b) = \{\{a\}, \{a, b\}\}$, rigorously prove that $(a, b) = (c, d)$ if and only if $a = c$ and $b = d$.
\end{problem}

\begin{problem}
Let $A = \{x\}$ and $B = \{y, z\}$. Explicitly enumerate the elements of $\mathcal{P}(A \times B)$.
\end{problem}

\begin{problem}
Let $A$ and $B$ be finite sets. Prove that $|\mathcal{P}(A \times B)| = 2^{|A| \cdot |B|}$.
\end{problem}

\begin{problem}
Find a counterexample to the claim: $\mathcal{P}(A \times B) = \mathcal{P}(A) \times \mathcal{P}(B)$. Explain the fundamental type mismatch occurring in this false claim.
\end{problem}

\begin{problem}
Prove that for any sets $A, B, C,$ and $D$, $\mathcal{P}(A \times B) \cap \mathcal{P}(C \times D) = \mathcal{P}((A \cap C) \times (B \cap D))$.
\end{problem}

\subsection{Empty Sets, Vacuous Truths, and Contrapositives}
\begin{problem}
Prove that $A \times B = \emptyset$ if and only if $A = \emptyset$ or $B = \emptyset$.
\end{problem}

\begin{problem}
Prove that if $A \times B \subseteq C \times D$ and $A \times B \neq \emptyset$, then $A \subseteq C$ and $B \subseteq D$.
\end{problem}

\begin{problem}
Provide a counterexample to the following statement: "If $A \times B = A \times C$, then $B = C$."
\end{problem}

\begin{problem}
Prove that if $A \times B = A \times C$ and $A \neq \emptyset$, then $B = C$.
\end{problem}

\begin{problem}
Prove that $(A \times B) \setminus (C \times D) = \emptyset$ if and only if $(A \subseteq C \text{ and } B \subseteq D)$ or $A = \emptyset$ or $B = \emptyset$.
\end{problem}

\subsection{Distributivity and Logical Equivalences}
\begin{problem}
Prove that $A \times (B \cup C) = (A \times B) \cup (A \times C)$ using logical equivalences.
\end{problem}

\begin{problem}
Prove that $A \times (B \cap C) = (A \times B) \cap (A \times C)$ using logical equivalences.
\end{problem}

\begin{problem}
Prove that $A \times (B \setminus C) = (A \times B) \setminus (A \times C)$.
\end{problem}

\begin{problem}
Prove that $(A \times B) \cap (C \times D) = (A \cap C) \times (B \cap D)$.
\end{problem}

\begin{problem}
Establish the logical equivalence of the following set difference of Cartesian products: 
\[ (A \times B) \setminus (C \times D) = ((A \setminus C) \times B) \cup (A \times (B \setminus D)) \]
\end{problem}

\subsection{Subset Inclusions, Unions, and Complements}
\begin{problem}
Prove by cases that $(A \times B) \cup (C \times D) \subseteq (A \cup C) \times (B \cup D)$.
\end{problem}

\begin{problem}
Provide a concrete counterexample showing that $(A \cup C) \times (B \cup D) \subseteq (A \times B) \cup (C \times D)$ is generally false.
\end{problem}

\begin{problem}
Let $U$ be a universal set such that $A \subseteq U$ and $B \subseteq U$. Let $S^c$ denote the complement of $S$ with respect to $U$. Prove that $(A \times B)^c = (A^c \times U) \cup (U \times B^c)$.
\end{problem}

\begin{problem}
Prove that $(A \times B) \cap (C \times D) = \emptyset$ if and only if $A \cap C = \emptyset$ or $B \cap D = \emptyset$.
\end{problem}

\begin{problem}
Prove that $(A \setminus B) \times (C \setminus D) \subseteq (A \times C) \setminus (B \times D)$.
\end{problem}

\subsection{Infinite Cartesian Products and Advanced Set Theory}
\begin{problem}
Let $I$ be an index set and $\{A_i\}_{i \in I}$ be a family of sets. Prove the contrapositive of the Axiom of Choice equivalence: If $\prod_{i \in I} A_i \neq \emptyset$, then $A_i \neq \emptyset$ for all $i \in I$.
\end{problem}

\begin{problem}
Let $\{A_i\}_{i \in I}$ and $\{B_i\}_{i \in I}$ be families of sets over the same index set $I$. Prove that $\prod_{i \in I} (A_i \cap B_i) = \left(\prod_{i \in I} A_i\right) \cap \left(\prod_{i \in I} B_i\right)$.
\end{problem}

\begin{problem}
Let $A_i = [0, 1/i]$ for $i \in \mathbb{N}^+$. Describe the infinite Cartesian product $\prod_{i \in \mathbb{N}^+} A_i$ and prove it contains exactly one element.
\end{problem}

\subsection{Abstract Real-World Modeling (Probability \& MDPs)}
\begin{problem}
\textbf{(Joint Probability Spaces)} Let $\Omega = X \times Y$ be a joint probability space. Assume the probability measure $\mathbb{P}$ factors such that $\mathbb{P}(A \times B) = \mathbb{P}_X(A)\mathbb{P}_Y(B)$ for all $A \subseteq X, B \subseteq Y$. Prove that $\mathbb{P}((A^c \times Y) \cap (X \times B^c)) = \mathbb{P}_X(A^c)\mathbb{P}_Y(B^c)$.
\end{problem}

\begin{problem}
\textbf{(State-Action Spaces)} In a Markov Decision Process, let the state space be $S = P \times V$ and the action space be $A = \{-1, 0, 1\}$. Define the "active trading" subset $\Omega_{active} \subseteq S \times A$ as the set of all state-action pairs where the action is non-zero. Prove that $\Omega_{active} = S \times \{-1, 1\}$.
\end{problem}

\begin{problem}
\textbf{(Bellman Domain)} Let $S$ and $A$ be finite sets representing states and actions. A deterministic policy is a function $\pi: S \to A$. Prove that the set of all possible deterministic policies is equivalent to the Cartesian product $\prod_{s \in S} A$, and determine its cardinality.
\end{problem}

\begin{problem}
\textbf{(Relations as Subsets)} A relation $R$ from $X$ to $Y$ is defined as $R \subseteq X \times Y$. The inverse relation is $R^{-1} = \{(y,x) \mid (x,y) \in R\}$. Prove that if $X=Y$, $R$ is symmetric if and only if $R = R^{-1}$.
\end{problem}

\subsection{Diagnostic Testing: "Find the Flaw"}
\begin{problem}
\textbf{Find the Flaw:} "Claim: If $A \times B \subseteq C \times D$, then $A \subseteq C$. Proof: Let $x \in A$. Choose $y \in B$. Then $(x,y) \in A \times B$. Since $A \times B \subseteq C \times D$, it follows that $(x,y) \in C \times D$. By definition, $x \in C$. Thus $A \subseteq C$."
\end{problem}

\begin{problem}
\textbf{Find the Flaw:} "Claim: $(A \cup B) \times (C \cup D) = (A \times C) \cup (B \times D)$. Proof: Let $(x,y) \in (A \cup B) \times (C \cup D)$. Then $x \in A \cup B$ and $y \in C \cup D$. This means $(x \in A \lor x \in B) \land (y \in C \lor y \in D)$. Distributing the AND over the OR gives $(x \in A \land y \in C) \lor (x \in B \land y \in D)$. Thus $(x,y) \in (A \times C) \cup (B \times D)$."
\end{problem}

\begin{problem}
\textbf{Find the Flaw:} "Claim: $\mathcal{P}(A \times B) = \{ (S_A, S_B) \mid S_A \subseteq A \text{ and } S_B \subseteq B \}$. Proof: Any subset of $A \times B$ must be formed by taking some subset of $A$ and pairing it with some subset of $B$. Therefore, the power set of the Cartesian product is just the set of all pairs of subsets."
\end{problem}

\newpage

% ==========================================
% SECTION 3: DETAILED SOLUTIONS
% ==========================================
\section{Detailed Solutions}
\label{sec:solutions}

\begin{proof}[Solution to Problem 1]
We must prove $(a,b) = (c,d) \iff a=c \land b=d$ using $(a,b) = \{\{a\}, \{a,b\}\}$. \\
\textbf{Forward Direction ($\implies$):} Assume $\{\{a\}, \{a,b\}\} = \{\{c\}, \{c,d\}\}$. By the axiom of extensionality, the elements of these sets must be equal. We proceed by cases based on whether $a=b$. \\
\textit{Case 1:} $a = b$. Then $(a,a) = \{\{a\}, \{a,a\}\} = \{\{a\}\}$. Thus, $\{\{c\}, \{c,d\}\} = \{\{a\}\}$. This implies $\{c\} = \{a\}$ and $\{c,d\} = \{a\}$. From $\{c\} = \{a\}$, we get $c=a$. From $\{c,d\} = \{a\}$, since $c=a$, we have $\{a,d\} = \{a\}$, which implies $d=a$. Thus $a=c$ and $b=a=d$, so $a=c \land b=d$. \\
\textit{Case 2:} $a \neq b$. Then $\{a\} \neq \{a,b\}$. Since the sets are equal, $\{c\}$ must equal one of the elements in $\{\{a\}, \{a,b\}\}$. Since $a \neq b$, $\{a,b\}$ has two elements, but $\{c\}$ has one. Thus $\{c\} \neq \{a,b\}$, forcing $\{c\} = \{a\}$, which means $c=a$. The remaining element must match: $\{c,d\} = \{a,b\}$. Since $c=a$, $\{a,d\} = \{a,b\}$. Since $a \neq b$, we must have $d=b$. Thus $a=c \land b=d$. \\
\textbf{Reverse Direction ($\impliedby$):} Assume $a=c$ and $b=d$. Substituting these into the Kuratowski definition yields $(a,b) = \{\{a\}, \{a,b\}\} = \{\{c\}, \{c,d\}\} = (c,d)$. 
\end{proof}

\begin{proof}[Solution to Problem 2]
First, compute the Cartesian product $A \times B$:
\[ A \times B = \{x\} \times \{y, z\} = \{(x, y), (x, z)\} \]
The power set $\mathcal{P}(A \times B)$ is the set of all subsets of $A \times B$. Since $|A \times B| = 2$, the power set will have $2^2 = 4$ elements:
\[ \mathcal{P}(A \times B) = \left\{ \emptyset, \{(x, y)\}, \{(x, z)\}, \{(x, y), (x, z)\} \right\} \]
\end{proof}

\begin{proof}[Solution to Problem 3]
By the fundamental counting principle, the number of elements in the Cartesian product $A \times B$ is the product of the number of elements in $A$ and $B$. Thus, $|A \times B| = |A| \cdot |B|$. 
The power set of any finite set $S$ has cardinality $2^{|S|}$. 
Substituting $S = A \times B$, we obtain:
\[ |\mathcal{P}(A \times B)| = 2^{|A \times B|} = 2^{|A| \cdot |B|} \]
\end{proof}

\begin{proof}[Solution to Problem 4]
Let $A = \{1\}$ and $B = \{2\}$. 
Then $A \times B = \{(1,2)\}$, and $\mathcal{P}(A \times B) = \{\emptyset, \{(1,2)\}\}$.
Next, $\mathcal{P}(A) = \{\emptyset, \{1\}\}$ and $\mathcal{P}(B) = \{\emptyset, \{2\}\}$.
The Cartesian product of the power sets is:
\[ \mathcal{P}(A) \times \mathcal{P}(B) = \{(\emptyset, \emptyset), (\emptyset, \{2\}), (\{1\}, \emptyset), (\{1\}, \{2\})\} \]
Clearly, $\mathcal{P}(A \times B) \neq \mathcal{P}(A) \times \mathcal{P}(B)$. 
\textbf{Type Mismatch:} $\mathcal{P}(A \times B)$ is a set containing sets of ordered pairs. $\mathcal{P}(A) \times \mathcal{P}(B)$ is a set containing ordered pairs of sets. These are fundamentally different mathematical objects.
\end{proof}

\begin{proof}[Solution to Problem 5]
We must show $\mathcal{P}(A \times B) \cap \mathcal{P}(C \times D) = \mathcal{P}((A \cap C) \times (B \cap D))$.
Let $S$ be an arbitrary set.
\begin{align*}
S \in \mathcal{P}(A \times B) \cap \mathcal{P}(C \times D) &\iff S \in \mathcal{P}(A \times B) \land S \in \mathcal{P}(C \times D) \\
&\iff S \subseteq (A \times B) \land S \subseteq (C \times D) \\
&\iff S \subseteq ((A \times B) \cap (C \times D)) \quad \text{(Definition of intersection of sets)} \\
&\iff S \subseteq ((A \cap C) \times (B \cap D)) \quad \text{(By Theorem 1.6 from notes)} \\
&\iff S \in \mathcal{P}((A \cap C) \times (B \cap D))
\end{align*}
Since the logical equivalences hold for all $S$, the sets are equal.
\end{proof}

\begin{proof}[Solution to Problem 6]
We prove $A \times B = \emptyset \iff A = \emptyset \lor B = \emptyset$. \\
\textbf{Forward ($\implies$):} We proceed by contrapositive. Assume $\neg(A = \emptyset \lor B = \emptyset)$, which by De Morgan's Law is $A \neq \emptyset \land B \neq \emptyset$. Since $A$ and $B$ are non-empty, there exists $a \in A$ and $b \in B$. By definition, the ordered pair $(a,b) \in A \times B$. Thus, $A \times B \neq \emptyset$. \\
\textbf{Reverse ($\impliedby$):} Assume $A = \emptyset \lor B = \emptyset$. Suppose for contradiction that $A \times B \neq \emptyset$. Then there exists some $(x,y) \in A \times B$, meaning $x \in A$ and $y \in B$. If $A = \emptyset$, no such $x$ exists (contradiction). If $B = \emptyset$, no such $y$ exists (contradiction). Thus, $A \times B = \emptyset$.
\end{proof}

\begin{proof}[Solution to Problem 7]
Assume $A \times B \subseteq C \times D$ and $A \times B \neq \emptyset$. 
Because $A \times B \neq \emptyset$, it follows from Problem 6 that $A \neq \emptyset$ and $B \neq \emptyset$.
To show $A \subseteq C$, let $x \in A$ be arbitrary. Since $B \neq \emptyset$, there exists some $y \in B$. Thus, $(x,y) \in A \times B$. Since $A \times B \subseteq C \times D$, $(x,y) \in C \times D$. By definition, $x \in C$ and $y \in D$. Thus $x \in C$, proving $A \subseteq C$.
To show $B \subseteq D$, let $y \in B$ be arbitrary. Since $A \neq \emptyset$, there exists some $x \in A$. Thus $(x,y) \in A \times B \implies (x,y) \in C \times D \implies y \in D$. Thus $B \subseteq D$.
\end{proof}

\begin{proof}[Solution to Problem 8]
Let $A = \emptyset$, $B = \{1\}$, and $C = \{2\}$. 
Then $A \times B = \emptyset \times \{1\} = \emptyset$.
And $A \times C = \emptyset \times \{2\} = \emptyset$.
Thus, $A \times B = A \times C$ holds (both are $\emptyset$). However, $B = \{1\} \neq \{2\} = C$. This serves as a valid counterexample.
\end{proof}

\begin{proof}[Solution to Problem 9]
Assume $A \times B = A \times C$ and $A \neq \emptyset$. We must show $B = C$.
To show $B \subseteq C$, let $y \in B$. Since $A \neq \emptyset$, there exists some $x \in A$. Thus, $(x,y) \in A \times B$. Since $A \times B = A \times C$, it follows that $(x,y) \in A \times C$. By definition of the Cartesian product, $y \in C$. Thus $B \subseteq C$.
By a symmetric argument, let $z \in C$. Since $x \in A$, $(x,z) \in A \times C = A \times B$, which implies $z \in B$. Thus $C \subseteq B$.
Therefore, $B = C$.
\end{proof}

\begin{proof}[Solution to Problem 10]
We know $S_1 \setminus S_2 = \emptyset \iff S_1 \subseteq S_2$. 
Thus, $(A \times B) \setminus (C \times D) = \emptyset \iff A \times B \subseteq C \times D$.
We proceed by cases on whether $A \times B$ is empty.
\textit{Case 1:} $A \times B = \emptyset$. By Problem 6, this occurs $\iff A = \emptyset \lor B = \emptyset$. In this case, $\emptyset \subseteq C \times D$ is vacuously true.
\textit{Case 2:} $A \times B \neq \emptyset$. This means $A \neq \emptyset \land B \neq \emptyset$. By Problem 7, $A \times B \subseteq C \times D \iff A \subseteq C \land B \subseteq D$.
Combining both cases with a logical OR, the condition holds if and only if $(A \subseteq C \land B \subseteq D) \lor A = \emptyset \lor B = \emptyset$.
\end{proof}

\begin{proof}[Solution to Problem 11]
Let $(x,y)$ be an arbitrary ordered pair.
\begin{align*}
(x,y) \in A \times (B \cup C) &\iff x \in A \land y \in (B \cup C) \\
&\iff x \in A \land (y \in B \lor y \in C) \\
&\iff (x \in A \land y \in B) \lor (x \in A \land y \in C) \quad \text{(Distributivity of } \land \text{ over } \lor) \\
&\iff (x,y) \in (A \times B) \lor (x,y) \in (A \times C) \\
&\iff (x,y) \in (A \times B) \cup (A \times C)
\end{align*}
Thus, $A \times (B \cup C) = (A \times B) \cup (A \times C)$.
\end{proof}

\begin{proof}[Solution to Problem 12]
Let $(x,y)$ be an arbitrary ordered pair.
\begin{align*}
(x,y) \in A \times (B \cap C) &\iff x \in A \land y \in (B \cap C) \\
&\iff x \in A \land (y \in B \land y \in C) \\
&\iff (x \in A \land y \in B) \land (x \in A \land y \in C) \quad \text{(Idempotence and Associativity of } \land) \\
&\iff (x,y) \in (A \times B) \land (x,y) \in (A \times C) \\
&\iff (x,y) \in (A \times B) \cap (A \times C)
\end{align*}
Thus, $A \times (B \cap C) = (A \times B) \cap (A \times C)$.
\end{proof}

\begin{proof}[Solution to Problem 13]
Let $(x,y)$ be an arbitrary ordered pair.
\begin{align*}
(x,y) \in A \times (B \setminus C) &\iff x \in A \land y \in (B \setminus C) \\
&\iff x \in A \land (y \in B \land y \notin C) \\
&\iff (x \in A \land y \in B) \land (x \in A \land y \notin C) \quad \text{(Since } x \in A \iff x \in A \land x \in A) \\
&\iff (x,y) \in (A \times B) \land (x,y) \notin (A \times C) \\
&\iff (x,y) \in (A \times B) \setminus (A \times C)
\end{align*}
Thus, $A \times (B \setminus C) = (A \times B) \setminus (A \times C)$.
\end{proof}

\begin{proof}[Solution to Problem 14]
Let $(x,y)$ be an arbitrary ordered pair.
\begin{align*}
(x,y) \in (A \times B) \cap (C \times D) &\iff (x,y) \in (A \times B) \land (x,y) \in (C \times D) \\
&\iff (x \in A \land y \in B) \land (x \in C \land y \in D) \\
&\iff (x \in A \land x \in C) \land (y \in B \land y \in D) \quad \text{(Commutativity of } \land) \\
&\iff x \in (A \cap C) \land y \in (B \cap D) \\
&\iff (x,y) \in (A \cap C) \times (B \cap D)
\end{align*}
Thus, $(A \times B) \cap (C \times D) = (A \cap C) \times (B \cap D)$.
\end{proof}

\begin{proof}[Solution to Problem 15]
Let $(x,y)$ be an arbitrary ordered pair.
\begin{align*}
(x,y) \in (A \times B) \setminus (C \times D) &\iff (x,y) \in (A \times B) \land (x,y) \notin (C \times D) \\
&\iff (x \in A \land y \in B) \land \neg(x \in C \land y \in D) \\
&\iff (x \in A \land y \in B) \land (x \notin C \lor y \notin D) \quad \text{(De Morgan's Law)} \\
&\iff (x \in A \land y \in B \land x \notin C) \lor (x \in A \land y \in B \land y \notin D) \quad \text{(Distributivity)} \\
&\iff ((x \in A \land x \notin C) \land y \in B) \lor (x \in A \land (y \in B \land y \notin D)) \\
&\iff (x \in (A \setminus C) \land y \in B) \lor (x \in A \land y \in (B \setminus D)) \\
&\iff (x,y) \in ((A \setminus C) \times B) \lor (x,y) \in (A \times (B \setminus D)) \\
&\iff (x,y) \in ((A \setminus C) \times B) \cup (A \times (B \setminus D))
\end{align*}
Thus, the sets are equal.
\end{proof}

\begin{proof}[Solution to Problem 16]
Let $(x,y) \in (A \times B) \cup (C \times D)$. By definition of union, $(x,y) \in A \times B$ or $(x,y) \in C \times D$. We proceed by cases. \\
\textit{Case 1:} $(x,y) \in A \times B$. Then $x \in A$ and $y \in B$. By the logical rule of addition, $x \in A \implies x \in A \lor x \in C \implies x \in A \cup C$. Similarly, $y \in B \implies y \in B \lor y \in D \implies y \in B \cup D$. Thus, $(x,y) \in (A \cup C) \times (B \cup D)$. \\
\textit{Case 2:} $(x,y) \in C \times D$. Then $x \in C$ and $y \in D$. By addition, $x \in C \implies x \in A \lor x \in C \implies x \in A \cup C$. Similarly, $y \in D \implies y \in B \lor y \in D \implies y \in B \cup D$. Thus, $(x,y) \in (A \cup C) \times (B \cup D)$. \\
In both cases, $(x,y) \in (A \cup C) \times (B \cup D)$, proving the subset inclusion.
\end{proof}

\begin{proof}[Solution to Problem 17]
Let $A = \{1\}$, $B = \{2\}$, $C = \{3\}$, and $D = \{4\}$.
Then $A \cup C = \{1, 3\}$ and $B \cup D = \{2, 4\}$.
The Cartesian product of the unions is:
\[ (A \cup C) \times (B \cup D) = \{(1,2), (1,4), (3,2), (3,4)\} \]
Now, compute the union of the Cartesian products:
\[ A \times B = \{(1,2)\} \quad \text{and} \quad C \times D = \{(3,4)\} \]
\[ (A \times B) \cup (C \times D) = \{(1,2), (3,4)\} \]
Notice that $(1,4) \in (A \cup C) \times (B \cup D)$, but $(1,4) \notin (A \times B) \cup (C \times D)$. Thus, the subset inclusion fails.
\end{proof}

\begin{proof}[Solution to Problem 18]
Let $(x,y) \in U \times U$ be an arbitrary ordered pair in the universal space.
\begin{align*}
(x,y) \in (A \times B)^c &\iff (x,y) \notin A \times B \\
&\iff \neg(x \in A \land y \in B) \\
&\iff x \notin A \lor y \notin B \quad \text{(De Morgan's Law)} \\
&\iff x \in A^c \lor y \in B^c
\end{align*}
Since $x, y \in U$, we can rewrite this as:
\begin{align*}
x \in A^c \lor y \in B^c &\iff (x \in A^c \land y \in U) \lor (x \in U \land y \in B^c) \\
&\iff (x,y) \in (A^c \times U) \lor (x,y) \in (U \times B^c) \\
&\iff (x,y) \in (A^c \times U) \cup (U \times B^c)
\end{align*}
Thus, $(A \times B)^c = (A^c \times U) \cup (U \times B^c)$.
\end{proof}

\begin{proof}[Solution to Problem 19]
By Theorem 1.6 (Problem 14), we know $(A \times B) \cap (C \times D) = (A \cap C) \times (B \cap D)$.
By the Empty Set Rule (Problem 6), a Cartesian product $X \times Y = \emptyset \iff X = \emptyset \lor Y = \emptyset$.
Substituting $X = A \cap C$ and $Y = B \cap D$, we get:
\[ (A \cap C) \times (B \cap D) = \emptyset \iff (A \cap C) = \emptyset \lor (B \cap D) = \emptyset \]
Thus, $(A \times B) \cap (C \times D) = \emptyset \iff A \cap C = \emptyset \lor B \cap D = \emptyset$.
\end{proof}

\begin{proof}[Solution to Problem 20]
Let $(x,y) \in (A \setminus B) \times (C \setminus D)$. 
By definition, $x \in A \setminus B$ and $y \in C \setminus D$.
This means $(x \in A \land x \notin B)$ and $(y \in C \land y \notin D)$.
Since $x \in A$ and $y \in C$, we have $(x,y) \in A \times C$.
We must show $(x,y) \notin B \times D$. Suppose for contradiction that $(x,y) \in B \times D$. Then $x \in B$ and $y \in D$. However, we already established $x \notin B$, which is a contradiction. Thus, $(x,y) \notin B \times D$.
Since $(x,y) \in A \times C$ and $(x,y) \notin B \times D$, it follows that $(x,y) \in (A \times C) \setminus (B \times D)$.
\end{proof}

\begin{proof}[Solution to Problem 21]
We must prove: If $\prod_{i \in I} A_i \neq \emptyset$, then $A_i \neq \emptyset$ for all $i \in I$.
Assume $\prod_{i \in I} A_i \neq \emptyset$. By definition of the infinite Cartesian product, there exists at least one function $f \in \prod_{i \in I} A_i$. 
By definition, $f: I \to \bigcup_{i \in I} A_i$ such that $f(i) \in A_i$ for every $i \in I$.
Suppose for contradiction that there exists some index $k \in I$ such that $A_k = \emptyset$. 
Then the function $f$ must map $k$ to some element $f(k) \in A_k$. But since $A_k = \emptyset$, no such element exists, meaning $f(k)$ is undefined. This contradicts the premise that $f$ is a valid function defined on all of $I$. 
Thus, no such empty $A_k$ can exist, meaning $A_i \neq \emptyset$ for all $i \in I$.
\end{proof}

\begin{proof}[Solution to Problem 22]
Let $f$ be an arbitrary function with domain $I$.
\begin{align*}
f \in \prod_{i \in I} (A_i \cap B_i) &\iff \forall i \in I, f(i) \in (A_i \cap B_i) \\
&\iff \forall i \in I, (f(i) \in A_i \land f(i) \in B_i) \\
&\iff (\forall i \in I, f(i) \in A_i) \land (\forall i \in I, f(i) \in B_i) \\
&\iff f \in \prod_{i \in I} A_i \land f \in \prod_{i \in I} B_i \\
&\iff f \in \left(\prod_{i \in I} A_i\right) \cap \left(\prod_{i \in I} B_i\right)
\end{align*}
Thus, the sets are equal.
\end{proof}

\begin{proof}[Solution to Problem 23]
We are given $A_i = [0, 1/i]$ for $i \in \mathbb{N}^+$. 
An element $f \in \prod_{i \in \mathbb{N}^+} A_i$ is a sequence $(x_1, x_2, x_3, \dots)$ such that $x_i \in [0, 1/i]$ for all $i \ge 1$.
Notice that as $i \to \infty$, the interval $[0, 1/i]$ shrinks. 
However, the Cartesian product is the set of all valid choice functions. For any $i$, $0 \in [0, 1/i]$. Thus, the sequence $(0, 0, 0, \dots)$ is a valid element, so the product is non-empty.
Wait, the problem asks to prove it contains exactly one element? No, $x_1$ can be $0.5 \in [0,1]$, $x_2$ can be $0.2 \in [0, 1/2]$. There are uncountably many sequences. 
\textit{Correction to Problem Statement interpretation:} The problem asks to describe it. It does NOT contain exactly one element. It contains uncountably many elements because $A_1 = [0,1]$ has uncountably many choices, $A_2$ has uncountably many, etc. 
Let us prove it is uncountably infinite. Since $A_1 = [0,1]$, we can define a family of functions $f_c$ for each $c \in [0,1]$ where $f_c(1) = c$ and $f_c(i) = 0$ for $i > 1$. Since there are uncountably many $c \in [0,1]$, there are uncountably many such functions. Thus the product is uncountably infinite.
\end{proof}

\begin{proof}[Solution to Problem 24]
We are given $\mathbb{P}(A \times B) = \mathbb{P}_X(A)\mathbb{P}_Y(B)$. We must find $\mathbb{P}((A^c \times Y) \cap (X \times B^c))$.
First, simplify the set intersection:
\[ (A^c \times Y) \cap (X \times B^c) = (A^c \cap X) \times (Y \cap B^c) = A^c \times B^c \]
Now, apply the probability measure:
\[ \mathbb{P}(A^c \times B^c) = \mathbb{P}_X(A^c)\mathbb{P}_Y(B^c) \]
This follows directly from the independence assumption of the joint probability space.
\end{proof}

\begin{proof}[Solution to Problem 25]
The state space is $S = P \times V$ and action space is $A = \{-1, 0, 1\}$.
The active trading subset is $\Omega_{active} = \{(s, a) \in S \times A \mid a \neq 0\}$.
By definition, $a \in \{-1, 0, 1\}$ and $a \neq 0 \implies a \in \{-1, 1\}$.
Thus, an element $(s,a) \in \Omega_{active} \iff s \in S \land a \in \{-1, 1\}$.
By definition of the Cartesian product, this is exactly the set $S \times \{-1, 1\}$.
\end{proof}

\begin{proof}[Solution to Problem 26]
A deterministic policy $\pi: S \to A$ assigns exactly one action $a \in A$ to each state $s \in S$.
By Definition 7.1 (Infinite/Arbitrary Cartesian Product), $\prod_{s \in S} A$ is the set of all functions $f: S \to A$. Thus, the set of all policies is exactly $\prod_{s \in S} A$.
To find the cardinality, since $S$ and $A$ are finite, for each of the $|S|$ states, there are $|A|$ possible independent choices of actions. 
Thus, the total number of functions is $|A| \times |A| \times \dots \times |A|$ ($|S|$ times), which equals $|A|^{|S|}$.
\end{proof}

\begin{proof}[Solution to Problem 27]
Assume $R \subseteq X \times X$. 
\textbf{Forward ($\implies$):} Assume $R$ is symmetric. By definition, $(x,y) \in R \implies (y,x) \in R$. 
Let $(a,b) \in R$. Then $(b,a) \in R$. By definition of inverse, $(a,b) \in R^{-1}$. Thus $R \subseteq R^{-1}$.
Let $(c,d) \in R^{-1}$. By definition of inverse, $(d,c) \in R$. Since $R$ is symmetric, $(c,d) \in R$. Thus $R^{-1} \subseteq R$. Therefore $R = R^{-1}$. \\
\textbf{Reverse ($\impliedby$):} Assume $R = R^{-1}$. Let $(x,y) \in R$. Since $R = R^{-1}$, $(x,y) \in R^{-1}$. By definition of inverse, this implies $(y,x) \in R$. Thus $R$ is symmetric.
\end{proof}

\begin{proof}[Solution to Problem 28]
\textbf{Flaw:} The proof assumes "Choose $y \in B$" is always a valid logical step. However, if $B = \emptyset$, there is no such $y \in B$ to choose. If $B = \emptyset$, then $A \times B = \emptyset$. The premise $\emptyset \subseteq C \times D$ is vacuously true regardless of what $A$ is. Therefore, $A$ does not have to be a subset of $C$. 
\textbf{Correction:} The theorem is false unless we explicitly add the hypothesis that $B \neq \emptyset$.
\end{proof}

\begin{proof}[Solution to Problem 29]
\textbf{Flaw:} The logical distribution step is entirely false. The proof claims that $(x \in A \lor x \in B) \land (y \in C \lor y \in D)$ is logically equivalent to $(x \in A \land y \in C) \lor (x \in B \land y \in D)$. 
In propositional logic, $(P \lor Q) \land (R \lor S)$ expands via FOIL to $(P \land R) \lor (P \land S) \lor (Q \land R) \lor (Q \land S)$. 
The proof arbitrarily dropped the "cross terms" $(x \in A \land y \in D)$ and $(x \in B \land y \in C)$. 
\textbf{Correction:} The correct expansion yields $(A \times C) \cup (A \times D) \cup (B \times C) \cup (B \times D)$.
\end{proof}

\begin{proof}[Solution to Problem 30]
\textbf{Flaw:} The proof conflates a subset of ordered pairs with an ordered pair of subsets. 
$\mathcal{P}(A \times B)$ is a set of subsets. An element $S \in \mathcal{P}(A \times B)$ is a set of ordered pairs, e.g., $S = \{(x_1, y_1), (x_2, y_2)\}$. 
The set $\{ (S_A, S_B) \mid S_A \subseteq A, S_B \subseteq B \}$ is actually $\mathcal{P}(A) \times \mathcal{P}(B)$. 
As proven in Problem 4, these two sets are fundamentally different data structures and have different cardinalities ($2^{|A||B|}$ vs $2^{|A|+|B|}$). 
\textbf{Correction:} The claim is mathematically nonsensical due to a severe type mismatch.
\end{proof}

\newpage

% ==========================================
% SECTION 4: CONCEPT CHECKLIST (REPEATED)
% ==========================================
\section{Concept Checklist \& Question Mapping (Repeated)}
\label{sec:checklist_end}

\begin{itemize}
    \item \textbf{C1: Kuratowski Definition of Ordered Pairs} -- Formal set-theoretic definition of $(a,b)$. (Tested in P1)
    \item \textbf{C2: Power Sets \& Combinatorial Enumeration} -- Cardinality and power sets of Cartesian products. (Tested in P2, P3, P4, P5)
    \item \textbf{C3: Empty Sets \& Vacuous Truths} -- Handling $\emptyset$ as the "zero" of set theory; existential instantiation. (Tested in P6, P7, P8, P9, P10)
    \item \textbf{C4: Distributivity of Cartesian Products} -- Distributing over $\cup, \cap, \setminus$. (Tested in P11, P12, P13, P14, P15)
    \item \textbf{C5: Subset Inclusion \& Proof by Cases} -- Handling unions and logical "OR" statements. (Tested in P16, P17, P20)
    \item \textbf{C6: Set Equality via Logical Equivalence} -- Translating set operations to propositional logic ($\land, \lor$). (Tested in P14, P15, P18, P19)
    \item \textbf{C7: Complements of Cartesian Products} -- De Morgan's Laws applied to 2D spaces. (Tested in P18)
    \item \textbf{C8: Infinite Cartesian Products \& Axiom of Choice} -- Arbitrary intersections and non-empty families of sets. (Tested in P21, P22, P23)
    \item \textbf{C9: Abstract Real-World Modeling} -- Markov Decision Processes (MDPs), State-Action spaces, and Joint Probability. (Tested in P24, P25, P26, P27)
    \item \textbf{C10: Diagnostic Testing (Find the Flaw)} -- Identifying logical fallacies, type mismatches, and false assumptions. (Tested in P28, P29, P30)
\end{itemize}

\end{document}